\scanpage{144}
% archive foliation: 143
une base de $\mathcal H\otimes_{\mathbb Z}\mathbb R$ (\inserted{fibre par fibre} sur $\mathbb R$), la donnée d'une telle structure complexe équivaut à la donnée de
\begin{equation}
\tau=z_2/z_1
\tag{3}
\end{equation}
en tant que section du faisceau structural $\mathcal O_S$. (\uline{NB} \deleted{il faut,} Pour une structure complexe donnée sur $\mathcal H\otimes_{\mathbb Z}\mathbb R$, $z_1$ et $z_2$ sont deux sections\inserted{-bases} du faisceau inversible obtenu, et $\tau$ est la section de $\underline{\mathcal O}_S$ qui fait passer de l'une à l'autre:
\begin{equation}
z_2=\tau z_1.
\tag{3 bis}
\end{equation}
\margin{Noter que $\forall s$, $\tau(s)\in\mathbb C\setminus\mathbb R$}

Une structure \uline{de Torelli} sur une \deleted{\unsure{fam}} courbe elliptique relative $E$ sur $S$ (correspondant à: $\mathcal H$ sur $S$ $+$ structure complexe sur $\mathcal H\otimes_{\mathbb Z}\mathbb R$) est la donnée d'un isomorphisme
\begin{equation}
t:\mathbb Z_S^2\xrightarrow{\sim}\mathcal H,
\tag{4}
\end{equation}
\deleted{\unsure{ou d'}}\inserted{ou encore} \deleted{une} de deux sections $z_1,z_2$ de $\mathcal H$, \deleted{\unsure{satisfais}} formant une base de $\mathcal H$ en tte fibre. La catégorie des courbes elliptiques relatives, sur $S$, avec structure de Torelli, \unsure{c'est|est} \deleted{équivalente à}\inserted{``rigide''}\deleted{est} (pas d'autom. $\ne\mathrm{id}$ d'un objet) et l'ens des
